


\section*{Phaseability}

% define min. read len as fn of heterozygosity
% => handle polyploids: multiply heterozygous vars?
% => account for pop structure/IBD afterwards
% consrtuct single probability output, then solve for read len or heterozygosity

In this section, we want to obtain minimal input requirements for reconstructing phased potato genomes, and look at how different properties of our sequencing technology and genome characteristics affect these requirements.
To approach this problem, some clarity on what phasing is and is not is required, in particular how it relates to genome assembly.
In the cells of a given biological organism, there are often multiple copies of the genome – depending on species and tissue, 10 or more are not unheard of.
Because of standing variation or spontaneous mutational processes, these so-called haplotypes may differ between another in minor or major ways.
Variation on the same haplotype is inherited together, unless recombination breaks it up; haplotypes are thus also referred to as the units of heredity.
We term a position at which a haplotype is different from all others a haplo-unique site; they may correspond to single small variants, (parts of) larger variants, or clusters of multiple nearby variants.
A phasing of a genome refers to a map taking two different haplo-unique sites, and telling us whether they are found on the same haplotype or not.
How do we construct such a phasing?
We will mainly concern ourselves with the minimal requirements for constructing a phasing from a set of genomic sequencing reads, but it should be noted that HiC sequencing, a mapping population or parental sequencing may also serve to accomplish this goal, though less economically. % justify why only deriving minimal estimates/can assemblers reach this?
When encountering a read, we know all of its sequence to be derived from a single physical piece of DNA corresponding to a fragment of a single haplotype\footnote{With the possible exception of PCR artefacts.}.
If two haplo-unique sequences $u_0, u_1$ are found on the same read $(u_0, u_1)$, we can assign them to the same haplotype; this construction can be transitively extended, if we have two reads $(u_0, u_1), (u_1, u_2)$, we can use the shared haplo-unique sequence $u_1$ to assign $u_0, u_2$ to the same haplotype.
Note that, contrary to genome assembly, repetitive sequences are not inherently problematic; so long as the sequence contains sufficient haplo-unique sequences, it can be easy to phase but difficult to assemble; vice versa, a perfectly nonrepetitive sequence may be difficult or impossible to phase but easy to assemble if it contains large stretches shared by all haplotypes.

We can thus derive a simple minimal requirement for reconstructing the phasing of a genome, given in theorem \ref{theo-phase-simple}

\begin{theorem}
	Given a set of reads $\mathcal{R}$ and a set of haplotypes $\mathcal{H}$, each modelled as consisting of an ordered set of sequences unique to it $\mathcal{H}_i := (u^i_0, u^i_1, \dots u^i_n)$, it is possible to construct a phasing $\mathcal{P}$ iff for every adjacent pair of haplo-unique sequences $u_j, u_{j+1}$ there is a read covering both: $\forall (u_j, u_{j+1}) \exists x \in \mathcal{R}: (u_j, u_{j+1}) \in x$.

\label{theo-phase-simple}
\end{theorem}

%TODO proof
A proof is simple to construct: if two haplo-unique sites $u_i, u_j$ with $i < j$ are on the same haplotype, a chain of neighbouring haplo-unique sequences $u_i, u_{i+1}, \dots, u_{j-1}, u_{j}$ can be found and thus $\mathcal{P}(u_i, u_j)$ evaluates to true.
If we know all neighbouring haplo-unique sequences to be covered, the absence of such a chain implies that $u_i$ and $u_j$ do not share a haplotype, and thus $\matcal{P}(u_i, u_j)$ can safely evaluate to false. 
Conversely, if there exists a pair of haplounique sequences $u_i, u_{i+1}$ on the same haplotype that are not covered by any read, we will not be able to resolve them unambiguously to recover their phasing.

%TODO statmodelling
Using some statistical modelling, we can turn this simple minimal requirement into an estimate of the sequencing coverage required to confidently resolve the phasing of a genome.
To ensure full phasing, the requirement of theorem \ref{theo-phase-simple} has to be met everywhere along the genome.
In particular, this means that
\begin{enumerate}
	\item every haplotype has to be sampled
	\item every position along the genome has to be sampled, and
	\item the read covering every position on every haplotype has to actually cover both adjacent haplo-unique sequences.
\end{enumerate}
We can model each of these requirements in turn, allowing us to derive a sequencing coverage estimate.

To model the first requirement, if we consider a given position along the genome, every read mapping there will be drawn from all possible haplotypes with equal probability.
Sampling is with replacement, if we assume our sequencing library to have sufficient size.
Asking how many samples need to be drawn to cover every haplotype with confidence is equivalent to the Coupon Collector Problem, for which multiple approximations have been previously derived. % find original citation, derive

The second requirement hearkens back to the early days of estimating sequencing requirements for whole-genome shotgun sequences -- except instead of ensuring every position is covered at least once, we need every position to be covered in sufficient depth to make sure the first requirement is met.
If we model observed coverage along the genome given a specified mean coverage using a Poisson distribution, we can derive a requirement for the mean coverage by solving for the quantile of the Poisson distribution.
This corresponds exactly to the inverse of the upper $\Gamma$ function, a numeric approximation of which is available in the \texttt{zipfR} package.

Finally, not all reads covering a certain position may be informative, if they happen to cover just one of the haplo-unique sequences.
We can compute the fraction of informative reads spanning both haplo-unique sequences from the fraction of valid starting positions, given that every read is uniquely specified by its starting position.
A starting position is valid if a read of length $\ell$ starting from it spans both haplo-unique sequences; all valid starting positions therefore have to be before the first haplo-unique sequences, and there will be $\ell-d$ many, with $d$ giving the distance between the two haplo-unique sequences.
As every read has a single starting position, a mean sequencing coverage of $X$ coverage corresponds to $X/\ell$ reads starting at every position along the genome in expectation.
We thus expect $(\ell-d)\frac{X}{\ell}=\frac{\ell-d}{\ell}X$ reads to fulfil this criterion.
As we assume reads to be Poisson distributed, we can simply multiply our coverage estimate derived in step 2 with this factor.

% incorporate after notion of msyd graph construction

\section*{Proof of Non-Circularity}
\newcommand{\aln}{\mathcal{A}}

We consider two multisyntenic regions $M$ and $M'$ that are both upstream and downstream of each other in genomes $q$ and $q'$, forming a circle in the graph representation.
Because downstream/upstreamness is transitive, this case covers all possible circles in the graph.
Formally, we can express this as $\exists q, q': M[q] < M'[q] \wedge M'[q'] < M[q']$.
Consider the alignment $\aln_{q \to q'}$.
If this hasn't explicitly been constructed, we can induce it from two alignments to a common reference $r$: $\forall i \aln_{q \to q'}:= \aln^{-1}_{q' \to r}[\aln_{q \to r}[i]]$.
Such an $r$ must exist, and be either the reference used during the realignment leading to $M$ or $M'$.
We can quickly see the two parts of the statement leading to a contradiction when we consider that multisyntenic regions are defined to be aligning on all organisms they cover:% at $\aln_{q \to q'}$ , we can see a contradiction emerge quickly to our :

\begin{eqnarray*}
M[q] &< M'[q] &\\
\aln_{q\to q'}[M[q]] &< \aln_{q \to q'}[M'[q]] &\mathrm{Alignments are monotonic}\\
M[q'] &< M'[q'] &\mathrm{but} M'[q'] < M'[q]\mathrm{!}
\end{eqnarray*}

%There are a few assumptions underlying this proof, some of which may be violated in practice.
Note that this argument assumes that our realignment finds monotonic and correct alignments.
\texttt{minimap2}'s fast approximative minimizer-based algorithm is used for computing alignments, which may in practice sometimes violate these conditions.
An additional source of alignment ambiguity is the stitching together of leftover non-multisyntenic sequence during the realignment step.\\
This proof does thus not guarantee that an msyd run may not produce output inducing a circular graph representation in certain edgecases, but instead serve to justify its pruning as making an approximative component of the algorithm match the mathematical guarantees of an exact algorithm.
